\chapter{Local Calculations and Nonlinear Constructions}
\label{ch:local-models}

A named function should not have to repeat the foundations every time it
is differentiated, composed, or evaluated in a new region.  What can be
reused is not an unrestricted function space but a finite description
with error bounds.  This chapter develops that description for local
series, and a second construction for nonlinear inverses and evolutions.
The earlier convex derivative remains the first route to the FTC; the
present local estimates provide further concrete FTC certificates.

\section{A series chart and its finite jets}
A chart about a rational complex point $c$ supplies coefficient
computations $a_n$ and positive rational constants $M,R$ with
\[
 F(c+w)=\sum_{n\ge0}a_nw^n,\qquad |a_n|\le M R^{-n}.
\]
The displayed equality names the prefix-and-tail computation, not an
independently assumed function.  The disk means inputs with a certified
bound $|w|\le r<R$.  Complex absolute values are computed by square roots;
all required upper bounds can be rational.  Put $q=r/R$.

A finite jet of order $k$ is the list of the first $k$ derivatives at a
specified center.  Coefficients give that list by $F^{(j)}(c)=j!a_j$.
For evaluation away from the center, the derivative of order $k$ has the
explicit tail, after original coefficient $N\ge k-1$,
\[
 T_{N,k}=\frac{M}{R^k}\sum_{n>N}(n)_k q^{n-k},
 \qquad (n)_k=n(n-1)\cdots(n-k+1).
\]
Here and below $(n)_k$ in a derivative-tail formula is a falling factorial;
rising factorials will always be identified when introduced.  This tail
is computable without an infinite summation procedure.  Once
\[
 q\frac{n+1}{n+1-k}\le\theta<1,
\]
the successive terms decrease geometrically; their remaining sum is at
most the first one divided by $1-\theta$.  The ratio decreases with $n$,
so one finite threshold proves the estimate for every later term.  The
case $q=0$ is direct coefficient evaluation.  Coefficient-box errors are
allocated separately to the finitely many retained terms.

The full derivative bound is particularly simple:
\[
 |F^{(k)}(c+w)|\le\frac{M k!}{R^k(1-q)^{k+1}}.
\]
It follows by repeatedly differentiating the finite geometric identity
and bounding its terminal terms.  Thus the chart supplies derivatives of
any requested finite order, rather than asking a value algorithm to reveal
its derivatives.

\section{Moving the center}
\begin{theorem}[Recentring with a new majorant]
\label{thm:recenter-model}
Let $|d|\le d_0<R$.  About $c+d$, the same computation has coefficients
\[
 b_k=\sum_{n=k}^{\infty}\binom nk a_nd^{n-k},
 \qquad |b_k|\le\frac{MR}{(R-d_0)^{k+1}}.
\]
They form a chart with radius $R-d_0$ and constant $MR/(R-d_0)$.
\end{theorem}
\begin{proof}
For a finite polynomial the change of center is the binomial identity.
The majorant for $b_k$ is
\[
 \frac{M}{R^k}\sum_{j\ge0}\binom{k+j}{k}(d_0/R)^j
 =\frac{M}{R^k}(1-d_0/R)^{-k-1}.
\]
The equality follows from the differentiated geometric identity.  To
compute $b_k$, bound the tail by the eventual ratio
$(d_0/R)(k+j+1)/(j+1)<1$.  On $|v|+d_0<R$, the omitted double sum is
absolutely bounded by a geometric series.  Finite binomial rearrangements
therefore prove $F(c+d+v)=\sum b_kv^k$ with explicit errors.
\end{proof}

Recentring alone does not create continuation beyond a certified original
domain: its new disk is contained in that domain.  To move farther one
needs new bounds supplied by an equation, an integral, or another proved
representation.  This distinction prevents an apparent continuation
algorithm from repeatedly spending a radius it never replenishes.

\section{Arithmetic, reciprocals, and composition}
Finite polynomial arithmetic and the same majorants give sum and product
computations.  For instance, two bounds $M R^{-n}$ and $N R^{-n}$ give
product coefficients at most $MN(n+1)R^{-n}$.  At any smaller radius
$S<R$, these are at most
\[
 \frac{MN}{(1-S/R)^2}S^{-n}.
\]
Restriction costs a bound but gives another chart of the same kind.

For a reciprocal write $F=a_0+H$, supply $|a_0|\ge m>0$, and choose $r$
so that the coefficient majorant for $H$ on that disk is at most
$\theta m$, with $\theta<1$.  Then
\[
 \frac1F=\frac1{a_0}\sum_{j\ge0}(-H/a_0)^j,
 \qquad
 \left|\text{tail after }j=N\right|
 \le\frac{\theta^{N+1}}{m(1-\theta)}.
\]
The coefficients can instead be generated by the finite recurrence
\[
 b_0=a_0^{-1},\qquad
 b_n=-a_0^{-1}\sum_{j=1}^na_jb_{n-j}.
\]
The geometric error proves agreement.  There is no decision of whether an
arbitrary computed $a_0$ is zero; the separation is an explicit input.

For composition, recenter the outer chart at $F(c)$ and write
$F(c+w)=F(c)+H(w)$, where $H$ has zero constant coefficient.  Choose an
inner radius on which its absolute coefficient majorant is at most
$s<R_G$, the outer radius.  The first $N$ coefficients of
$G(F(c)+H)$ require only finitely many coefficients of $H$, since
$H^j$ has order at least $j$.  The outer tail is bounded by
\[
 M_G(s/R_G)^{N+1}/(1-s/R_G).
\]
The absolute sum of all composed coefficients weighted by the chosen inner
radius is at most $M_G/(1-s/R_G)$.  Hence each coefficient has that
geometric majorant, and degree truncation on a strictly smaller disk has
its own geometric tail, separate from the outer-series truncation.
Finite errors in evaluating $H$ are multiplied by the outer first-derivative
bound on a slightly larger certified disk.  This proves both the
composition computation and its chain rule.  Computed centers cause no
new problem: all their uses are interval operations with the displayed
strict margins.

These constructions also resolve removable expressions.  For example,
instead of dividing $S(x)$ by $x$ at zero, shift the coefficients of its
odd series.  This defines $S(x)/x$ there with value $\pi$.  A cancellation
proved in the coefficient algebra is different from a numerical division
by a small, uncertain denominator.

\section{Which approximations support which operations?}
Suppose finite computations $F_N$ approximate a function uniformly to
error $e_N$ on $[a,b]$.  Then their integral computations differ by at most
$(b-a)e_N$.  This is a finite rectangle estimate and proves integration
of the limit computation whenever the $e_N$ shrink effectively.

To differentiate, more is needed.  If $F_N'$ converge uniformly to a
computed $G$, then their finite FTC identities give
\[
 F(y)-F(x)=\int_x^yG(t)\,dt.
\]
A supplied modulus for $G$ makes this a derivative and FTC certificate
for $F$.  Uniform convergence of values alone gives no such conclusion.
The derivative tails of a series chart supply exactly the stronger data.
This is the form of a convergence theorem that later constructions use.


\section{A smooth computation which is not a single analytic chart}
\label{sec:flat-cutoff}
Define at rational inputs
\[
 B(x)=\begin{cases}E(-1/x^2),&x\ne0,\\0,&x=0.\end{cases}
\]
Rational equality decides the displayed branch.  Away from zero its
successive derivatives have the explicit form
\[
 D^nB(x)=P_n(1/x)E(-1/x^2),\qquad
 P_0(u)=1,\quad P_{n+1}(u)=2u^3P_n(u)-u^2P_n'(u).
\]
This is a finite rational polynomial recurrence.  To control the missing
point, let $d$ be the degree of $P_n$ and choose an integer $m$ with
$2m\ge d+2$.  For $0<|x|\le\delta\le1$, positivity of the exponential
series gives $E(1/x^2)\ge |x|^{-2m}/m!$.  Writing
$P_n(u)=\sum_j p_ju^j$, it follows that
\[
 |P_n(1/x)E(-1/x^2)|
 \le m!\sum_j|p_j|\,|x|^{2m-j}
 \le m!\left(\sum_j|p_j|\right)\delta^2.
\]
In particular each derivative extends continuously by zero, and its
quotient by $x$ tends to zero with a linear error bound.  Induction proves
that these extensions are the actual successive derivatives, all zero
at the origin.  This is not an inference from pointwise differentiability
on a punctured rational interval; the estimate explicitly controls the join.

Bounds on a compact interval are obtained by splitting at rational
$\pm\delta$: the estimate handles the center, while ordinary series and
reciprocal bounds handle the two separated pieces.  The same procedure
supplies moduli for every specified derivative order.  The FTC across the
join follows from the certified identities on the separated pieces and
the vanishing bounded end errors.

For rational $a<b$, set $\psi(x)=B(x-a)B(b-x)$ on $a<x<b$ and zero
outside.  All its derivatives match zero at both endpoints by the finite
product rule and the preceding estimates.  This constructs a compactly
supported smooth test with any requested derivative and error bound.
The positive value at an interior rational point shows that it is not
identically zero, although every derivative at an endpoint vanishes.
No analytic chart spanning that endpoint can represent it.  Explicit
nonanalytic functions therefore belong to the same calculus, by their
own joining estimates rather than by an assertion of analyticity.

\section{A nonlinear inverse: Lambert's function}
\label{sec:lambert-construction}
For $|z|\le1/4$, start with $w_0=0$ and compute
\[
 w_{n+1}=zE(-w_n).
\]
The exponential series gives $E(1/2)<2$ by a finite prefix and tail.
Consequently this map sends $|w|\le1/2$ into itself and has Lipschitz
constant at most $1/2$ there.  The estimate follows either from its
uniform increment identity along the segment or directly from the
exponential difference formula.  Finite iteration proves
\[
 |w_m-w_n|\le 2^{1-n}|z|\qquad(m\ge n).
\]
Successively accurate interval evaluations, with a geometric error budget
for their finite compositions, give a valid computation $W(z)$.  It
satisfies $W E(W)=z$ by the iteration identity and continuity.  Two fixed
points in this disk are equal by repeated contraction.  This constructs
one specified branch; it does not silently choose a branch elsewhere.

Parameter comparison in the fixed-point equations gives
$|W(z)-W(v)|\le4|z-v|$.  The finite exponential remainder for
$P(w)=wE(w)$ has linear coefficient $E(w)(1+w)$, whose modulus on this
disk is at least $1/4$: $|1+w|\ge1/2$ and $|E(w)|\ge1/2$.
Substitute $w=W(z)$ and $w+\Delta w=W(z+h)$ into this remainder and use
$|\Delta w|\le4|h|$.  Division by the separated linear coefficient gives
\[
 DW(z)=\frac{E(-W(z))}{1+W(z)},
\]
with an explicit quadratic increment remainder.  All constants can be
read from the exponential bounds on $|w|\le1/2$.  On real subintervals,
that remainder also supplies the convex-difference FTC certificate by
adding a sufficiently large rational multiple of $x^2$.

The example is named in the standard notation of
\href{https://dlmf.nist.gov/4.13}{DLMF, Section 4.13}; the finite iteration
and its error estimate are the definition used here.

\section{Nonlinear initial-value problems on a certified slab}
\label{sec:nonlinear-slab}
Let a specific computation $Q(t,y)$ have interval-continuity data on
$|t-t_0|\le h$, $\|y-y_0\|\le r$, with rational bounds
\[
 \|Q\|\le B,\qquad
 \|Q(t,y)-Q(t,v)\|\le L\|y-v\|,\qquad
 hB\le r,\quad q=hL<1.
\]
The vector norm is the maximum norm.  These are bounds to be proved for
the displayed formula $Q$, not consequences of a general ODE existence
theorem.  Define finite Picard computations
\[
 y_0(t)=y_0,\qquad
 y_{n+1}(t)=y_0+\int_{t_0}^t Q(u,y_n(u))\,du.
\]
Every iterate stays in the box, and finite induction gives
\[
 \|y_{n+1}-y_n\|\le hBq^n,\qquad
 \|y_m-y_n\|\le\frac{hBq^n}{1-q}\quad(m\ge n).
\]
Thus the iterates construct a uniformly approximated solution.  Passing
through the bounded interval integrals gives its integral equation;
accumulation then proves $Dy=Q(t,y)$.  Two supplied integral solutions
in the slab agree by the same contraction estimate.

For example, $y'=1+y^2$, $y(0)=0$ has $r=1$, $h=1/4$, $B=2$,
$L=2$, and $q=1/2$.  Its finite iterates are rational polynomials.
On this interval the already constructed computation
$S(t/\pi)/C(t/\pi)$ has the same equation and initial value; the cosine
is positively separated and its increment identity proves the integral
equation.  The two computations therefore agree.  Continue only through
further certified slabs; no continuation through the tangent's pole is
asserted.  Nonlinear special-function equations can be admitted in this
same concrete manner without forcing them into the linear theory.
